\documentclass[a4paper, 14pt]{extarticle}

\usepackage{fontspec}
\usepackage{polyglossia}
\setdefaultlanguage{russian}
\setotherlanguages{english}

\newfontfamily\cyrillicfont{Times New Roman}
\setmainfont{Times New Roman}

\usepackage{amssymb}
\usepackage{amsmath}

\title{Сравнительный анализ численных методов в теории дифференциальных уравнений}
\author{Халгаев Давид}

\begin{document}
\maketitle

\section{Введение}
Обыкновенные дифференциальные уравнения (ОДУ), в отличие от уравнений в частных производных, — это 
соотношения между функциями и их производными, зависящими только от одной независимой переменной. 

\paragraph{}
В общетеоретическом плане ОДУ давно и достаточно хорошо исследованы, особенно для линейных 
случаев. Первые дифференциальные уравнения были сформулированы и решены изобретателями 
дифференциального и интегрального исчисления Исааком Ньютоном \\ (1642–1727 гг.) и Готфридом 
Лейбницем (1646–1716 гг.). Начала теории обыкновенных дифференциальных уравнений порождены еще 
братьями Бернулли Якобом (1654–1705 гг.)
и Иоганном (1667–1748 гг.), Л.Эйлером (1707–1783 гг.) — автором
метода ломаных, первого численного алгоритма решения ОДУ,
Ж.Лагранжем (1736–1813 гг.), за которым Л.Эйлер признал открытие вариационного исчисления, К. 
Гауссом \\ (1777–1855 гг.), О.Коши (1836 г. — первое доказательство теоремы существования решения 
системы ОДУ) и основоположником теории устойчивости А.М.Ляпуновым (1857–1918 гг.).

\section{Определение обыкновенного дифференциального уравнения и задача Коши}
Обыкновенным дифференциальным уравнением называет соотношение вида
$$ F(t,y(t), y'(t),\ldots,y^{(n)}(t)) = 0, $$

где $F$ - известная функция, $t$ - независимая переменная, а $y(t)$ - неизвестная функция.
Максимальный порядок $n$ производной функции $y(t)$, входящий в Оду, называется
порядком уравнения.

\par 
Рассмотрим ОДУ первого порядка
\[y'=f(x,y)\]
с начальным условием
\[y(x_0)=y_0\]
Эта задача называется задачей Коши.

\section{Существование и единственность решения задачи Коши}

\par \textbf{Определение.}  Условие Липшица выполняется для функции \(f(x)\) если,
\[\exists L :\forall x_1,x_2 ~ |f(x_1)-f(x_2)|\leqslant L|x_1-x_2|\]

\par
Пусть \(y'=f(x,y)\), функция \(f(x,y)\) определена по первому аргументу на отрезке 
\([x_0,x_0+X]\), по второму аргументу - на всей числовой прямой; \(f(x,y)\) - непрерывна
в области определения функции; выполняется условие Липшица по второму аргументу, т.е.
\[\exists L:\forall y_1,y_2 ~ |f(x,y_1)-f(x,y_2)\leqslant L|y_1-y_2|\]
Тогда существует единственное решение \(y(x)\) задачи Коши [2, 4 c.]

\section{Численные методы решения ОДУ}
Пусть отрезок \([x_0,x_0+X]\) на \(N\) равных частей с шагом \(h=\frac{X}{N}\), обозначив 
узлы \(x_i=x_0+ih,~i=0,\ldots,N\). Приближенные значения решения в узлах будем обозначать \(y_i\),
точное решение \(y(x_i)\).
\par
Задача состоит в конструировании алгоритма для \(y_i\) и исследования условий, при которых 
точные и приближенные решения близки.
\par \textbf{Определение.} Назовем метод сходящимся с порядком \(p>0\), если найдется такое \(C\), что выполняется 
\[|y_n-y(x_n)|\leqslant Ch^p ~ \forall n =1,\ldots,N.\]
\par \textbf{Определение.} Невязкой метода называется разность точного значение и значения метода в узле.

\section{Семейство явных методов Рунге-Кутты второго порядка}
\[y_{n+1}=y_n+\sigma_1k_1+\sigma_2k_2\]
- метод семейства Рунге-Кутты второго порядка, где
\[
\begin{cases}
    k_1 = hf(x_n,y_n) \\ 
    k_2 = hf(x_n+ah,y_n+bk_1) 
\end{cases}
\]
Семейство содержит четыре параметра: \(\sigma_1,\sigma_2,k_1,k_2\).

\par \textbf{Теорема.} Если явный метод Рунге-Кутты имеет
невязку порядка \((p + 1)\), то он сходится с порядком \(p\).

\subsection{Явный метод Эйлера}
При \(\sigma_1=1, \sigma_2=0\) получается явный метод Эйлера
\[y_{n+1}=y_n+hf(x_n,y_n)\]
явный метод Эйлера сходится с первым порядком.[2, с. 7]

\subsection{Метод Эйлера с пересчетом}
\(\sigma_1=\frac12,\sigma_2=\frac12,a=1,b=1\)
\[  y_1=y_0+\frac{h}{2}(f(x_0,y_0)+f(x_1,y_0+hf(x_0,y_0)))  \]
Метод Эйлера с пересчетом имеет невязку третьего порядка и сходится со степенью 2. [2, c. 9]

\subsection{Метод Коши}
\(\sigma_1=0,\sigma_2=1,a=\frac12,b=\frac12\)
\[ y_{n+1}=y_n+hf(x_n+\frac{h}{2},y_n+\frac{h}{2}f(x_n,y_n)) \]
метод Коши сходится со вторым порядком.[2, c. 11]

\section{Общий класс метотодов Рунге-Кутты}
Явными методами Рунге-Кутты называются алгоритмы
\[y_{n+1}=y_n+\sum^m_{i+1}\sigma_ik_i\]
где
\[k_1=hf(x_n,y_n)\]
\[k_i=hf(x_n+a_ih,y_n+\sum^{i-1}_{j=1}b_{ij}k_j\]
Число \(m\) называется числом этапов
\par
Неявными методами Рунге-Кутты называются алгоритмы
\[y_{n+1}=y_n+\sum^m_{i=1}\sigma_ik_i\]
где
\[k_i=hf(x_n+a_ih,y_n+\sum^m_{j=1}b_{ij}k_j\]

Для того чтобы реализовать неявные методы, в отличие от явных, необходимо решать нелинейную систему
относительно \(k_i\), поэтому эти методы сложны в реализации. Однако эти методы обладают рядом преимуществ, в 
частоности, ими можно решать жесткие системы.

Коэффициенты метода удобно задавать в виде таблицы - матрицы Бутчера:

\begin{center}
\begin{tabular}{c|c c c c}
     $a_1$ & $b_{11}$ & $b_{12}$ & $\ldots$ & $b_{1m}$ \\
     $a_2$ & $b_{21}$ & $b_{22}$ & $\ldots$ & $b_{2m}$ \\
     $\ldots$ & $\ldots$ & $\ldots$ & $\ldots$ & $\ldots$ \\
     $a_m$ & $b_{m1}$ & $b_{m2}$ & $\ldots$ & $b_{mm}$ \\
     \hline
      & $\sigma_1$ & $\sigma_2$ & $\ldots$ & $\sigma_m$
\end{tabular}    
\end{center}

Для явных методов таблица Бутчера выглядит так:

\begin{center}
\begin{tabular}{c|c c c c}
     0 & 0 & 0 & 0 & 0 \\
     $a_2$ & $b_{21}$ & 0 & $\ldots$ & 0 \\
     $\ldots$ & $\ldots$ & $\ldots$ & $\ldots$ & $\ldots$ \\
     $a_m$ & $b_{m1}$ & $b_{m2}$ & $\ldots$ & 0 \\
     \hline
      & $\sigma_1$ & $\sigma_2$ & $\ldots$ & $\sigma_m$
\end{tabular}    
\end{center}

Например, для метода Эйлера с пересчетом таблица 
записывается в виде

\begin{center}
\begin{tabular}{c|c c}
    0 & 0 & 0 \\
    1 & 1 & 0 \\ \hline
      & \(\frac12\) & \(\frac12\)
\end{tabular}
\end{center}

\section{Многошаговые методы}
Методы Рунге-Кутты основаны на многократном пересчитывании на каждом шаге значений правой части дифференциальных 
уравнений. Методы Адамса основаны на запоминании и использовании на каждом шаге уже посчитанных ранее значений правой 
части дифференциальных уравнений.

\subsection{Явный метод Адамса}
\[y_{n+1}=y_n+\int^{x_{n+1}}_{x_n}L_{k-1}(s)ds,\]
где \(L_{k-1}(s)\) - интерполяционный многочлен Лагранжа, построенный по вычисленным значениям \(f_n=f(x_n,y_n)\),
\(f_{n-1}=f(x_{n-1},y_{n-1}), \ldots,f_{n-k+1}=f(x_{n-k+1},y_{n-k+1})\). Число \(k\) называется числом шагов.
\par
При \(k=3\) многочлен имеет вид
\[L_2(x)=f_n+\frac{f_n-f_{n-1}}{h}(x-x_n)+\frac{f_n-2f_{n-1}+f_{n-2}}{2h^2}(x-x_n)(x-x_{n-1})\]
Подставляя в интеграл, получаем трехшаговый метод Адамса
\[y_{n+1}=y_n+\frac{h}{12}(23f_n-16f_{n-1}+5f_{n-2})\]
который называется явным методом Адамса третьего порядка или методом Адамса-Башфорта.

\subsection{Неявные методы Адамса}
Неявными методами Адамса называются методы
\[y_{n+1}=y_n+\int^{x_{n+1}}_{x_n}L_{k}(s)ds\]
эти метода также называются экстраполяционными методами Адамса,в отличие от явных, которые называются 
интерполяционными
\par 
При \(k=2\) получакм метод 
\[y_{n+1}=y_n+\frac{h}{12}(5f_{n+1}+8f_n-f_{n-1})\]
этот метод называется неявным методом Адамса третьего порядка или методом Адамса-Мултона.

\paragraph{}
Сравнивая явные и неявные методы Адамса, можно отметить следующее:
\par
1. Недостаток неявных методов состоит в необходимости на
каждом шаге решать уравнение относительно неизвестной величины \( y_{n+1} \).
\par
2. Некоторое преимущество неявных методов состоит в точности: при одной и той же шаговости к неявные методы имеют
порядок сходимости \((k + 1)\), в отличие от явных, у которых порядок сходимости \(k\).

\subsection{Общий класс методов многошаговых методов}
Явные \(k\)-шаговые методы имеют вид
\[y_n=\sum^k_{i=1}\alpha_iy_{n-i}+h\sum^k_{i=1}\beta_if_{n-i}\]
где \(f_{n-i}=f(x_{n-i},y_{n-i})\) \par
Неявные \(k\)-шаговые методы имеют вид
\[y_n=\sum^k_{i=1}\alpha_iy_{n-i}+h\sum^k_{i=0}\beta_if_{n-i}\]

\section{Таблица сравнения методов по порядку сходимости}

\begin{center}
\begin{tabular}{| l | l |}
    \hline
     Название метода & Порядок сходимости \\ \hline
     Явный метод Эйлера & 1 \\
     Метод Эйлера с пересчетом & 2 \\
     Метод Коши & 2 \\
     Метод Адамса-Башфорта \(k\)-го порядка & k \\
     Метод Адамса-Мултона \(k\)-го порядка & k+1 \\
     \hline
\end{tabular}
\end{center}

\newpage
\section*{Список литературы}

1. Методы решения обыкновенных дифференциальных уравнений, В.П.Ильин 
\par
2. Численные методы, В.Г. Пименов, А.Б. Ложников 

\end{document}

